\begin{table*}[t]\centering\caption{Main comparison (mean $\pm$ std over 5 seeds). valid\_time in model time units (higher is better), stable\_frac and offline\_r2 higher is better, other columns lower is better. Bold = best, underline = second best within a task.}\label{tab:main}
\footnotesize\resizebox{\textwidth}{!}{\begin{tabular}{llccccccccc}
\toprule
Method & Task & valid\_time $\uparrow$ & rmse\_0p5 $\downarrow$ & stable\_frac $\uparrow$ & offline\_r2 $\uparrow$ & mean\_err $\downarrow$ & var\_err $\downarrow$ & w1\_pdf $\downarrow$ & spec\_err $\downarrow$ & $n$ \\
\midrule
No closure & F20\_c10 & 0.363 $\pm$ 0.004 & 0.777 $\pm$ 0.019 & \textbf{1.000 $\pm$ 0.000} & 0.000 $\pm$ 0.000 & 0.439 $\pm$ 0.015 & 1.162 $\pm$ 0.008 & 1.895 $\pm$ 0.006 & 0.386 $\pm$ 0.002 & 5 \\
Polynomial (deg 4) & F20\_c10 & \underline{1.398 $\pm$ 0.041} & \underline{0.127 $\pm$ 0.007} & \underline{1.000 $\pm$ 0.000} & \underline{0.847 $\pm$ 0.001} & \underline{0.236 $\pm$ 0.038} & 0.042 $\pm$ 0.007 & \underline{0.243 $\pm$ 0.037} & 0.144 $\pm$ 0.020 & 5 \\
MLP (32x2) & F20\_c10 & \textbf{1.421 $\pm$ 0.036} & \textbf{0.120 $\pm$ 0.006} & 1.000 $\pm$ 0.000 & \textbf{0.850 $\pm$ 0.001} & 0.237 $\pm$ 0.046 & \underline{0.041 $\pm$ 0.008} & 0.245 $\pm$ 0.045 & \underline{0.137 $\pm$ 0.023} & 5 \\
\textbf{AR(1) stochastic} & F20\_c10 & 1.252 $\pm$ 0.029 & 0.145 $\pm$ 0.008 & 1.000 $\pm$ 0.000 & 0.847 $\pm$ 0.001 & \textbf{0.063 $\pm$ 0.007} & \textbf{0.002 $\pm$ 0.001} & \textbf{0.083 $\pm$ 0.003} & \textbf{0.056 $\pm$ 0.002} & 5 \\
\midrule
No closure & F20\_c4 & 0.673 $\pm$ 0.021 & 0.378 $\pm$ 0.030 & \textbf{1.000 $\pm$ 0.000} & 0.000 $\pm$ 0.000 & 0.252 $\pm$ 0.011 & 0.341 $\pm$ 0.003 & 0.851 $\pm$ 0.002 & 0.141 $\pm$ 0.001 & 5 \\
Polynomial (deg 4) & F20\_c4 & \underline{1.097 $\pm$ 0.035} & \underline{0.158 $\pm$ 0.021} & \underline{1.000 $\pm$ 0.000} & \underline{0.719 $\pm$ 0.003} & \textbf{0.007 $\pm$ 0.005} & \underline{0.003 $\pm$ 0.002} & \underline{0.030 $\pm$ 0.007} & 0.009 $\pm$ 0.002 & 5 \\
MLP (32x2) & F20\_c4 & \textbf{1.099 $\pm$ 0.031} & \textbf{0.157 $\pm$ 0.021} & 1.000 $\pm$ 0.000 & \textbf{0.720 $\pm$ 0.003} & \underline{0.009 $\pm$ 0.008} & \textbf{0.003 $\pm$ 0.002} & 0.032 $\pm$ 0.006 & \underline{0.008 $\pm$ 0.003} & 5 \\
\textbf{AR(1) stochastic} & F20\_c4 & 0.942 $\pm$ 0.031 & 0.179 $\pm$ 0.018 & 1.000 $\pm$ 0.000 & 0.719 $\pm$ 0.003 & 0.009 $\pm$ 0.006 & 0.005 $\pm$ 0.002 & \textbf{0.029 $\pm$ 0.005} & \textbf{0.007 $\pm$ 0.002} & 5 \\
\bottomrule
\end{tabular}
}
\end{table*}
\begin{table}[t]\centering\caption{Noise floor: an independent truth run scored against the reference (5 seeds).}\label{tab:floor}
\footnotesize\resizebox{\columnwidth}{!}{\begin{tabular}{llccccc}
\toprule
Method & Task & mean\_err $\downarrow$ & var\_err $\downarrow$ & w1\_pdf $\downarrow$ & spec\_err $\downarrow$ & $n$ \\
\midrule
Truth (indep. run) & F20\_c10 & \textbf{0.012 $\pm$ 0.004} & \textbf{0.002 $\pm$ 0.001} & \textbf{0.027 $\pm$ 0.005} & \textbf{0.013 $\pm$ 0.006} & 5 \\
\midrule
Truth (indep. run) & F20\_c4 & \textbf{0.020 $\pm$ 0.007} & \textbf{0.005 $\pm$ 0.001} & \textbf{0.031 $\pm$ 0.010} & \textbf{0.006 $\pm$ 0.001} & 5 \\
\bottomrule
\end{tabular}
}
\end{table}
\begin{figure}[t]\centering\includegraphics[width=0.85\columnwidth]{bars_main_valid_time.pdf}\caption{Forecast valid time by system and task (mean over 5 seeds; error bars: standard deviation over seeds).}\label{fig:valid}\end{figure}
\begin{figure}[t]\centering\includegraphics[width=\columnwidth]{spectrum_ratio.pdf}\caption{Log ratio of the time-mean spatial power at wavenumber $m$ to the reference (0 is perfect), mean $\pm$ std over 5 seeds.}\label{fig:spec}\end{figure}

\textbf{H1 (all learned closures beat no closure on valid time): supported.} In Table~\ref{tab:main}, at $c{=}10$ valid time rises from 0.363 (no closure) to 1.398 (polynomial), 1.421 (MLP) and 1.252 (AR(1)); at $c{=}4$ from 0.673 to 1.097, 1.099 and 0.942. All six differences have Welch $p<10^{-6}$ (n=5 each; \texttt{rh compare} with the no-closure reference). Normalised RMSE at lead 0.5 falls from 0.777 to 0.120--0.145 at $c{=}10$. This does not show skill beyond roughly one to one and a half time units: even the best closure has a mean valid time of 1.421, below the 2-unit forecast horizon.

\textbf{H2 (AR(1) has a better climate than the deterministic polynomial): supported at $c{=}10$, not at $c{=}4$; with a cost in skill.} At $c{=}10$ the deterministic closures give too little slow-variable variance (var\_err 0.042 for the polynomial against 0.002 for AR(1), $p=$\rhval{cmp/main/ar-1-stochastic/f20_c10/var_err/welch_p}) and a mis-shaped spectrum (spec\_err 0.144 against 0.056, $p=$\rhval{cmp/main/ar-1-stochastic/f20_c10/spec_err/welch_p}); the AR(1) PDF distance is also lower (w1\_pdf 0.083 against 0.243). Figure~\ref{fig:spec} shows the deterministic closures over-predict mode 1 and under-predict mode 4, which AR(1) reduces. AR(1) is still above the independent-truth floor in Table~\ref{tab:floor} for spec\_err (0.056 against 0.013) and w1\_pdf (0.083 against 0.027), but its var\_err (0.002) equals the floor. At $c{=}4$ the picture is different: spec\_err is 0.009 (polynomial), 0.008 (MLP), 0.007 (AR(1)), the difference polynomial--AR(1) has $p=$\rhval{cmp/main/ar-1-stochastic/f20_c4/spec_err/welch_p}, and var\_err is nominally worse for AR(1) (0.005 against 0.003, $p=$\rhval{cmp/main/ar-1-stochastic/f20_c4/var_err/welch_p}); all three are close to the floor (0.006 and 0.005), so the metrics cannot separate them. In both regimes AR(1) has a shorter valid time than the polynomial: 1.252 against 1.398 at $c{=}10$ ($p=$\rhval{cmp/main/ar-1-stochastic/f20_c10/valid_time/welch_p}) and 0.942 against 1.097 at $c{=}4$ ($p=$\rhval{cmp/main/ar-1-stochastic/f20_c4/valid_time/welch_p}; \texttt{rh compare} with the polynomial as reference). Figure~\ref{fig:tradeoff} places the $c{=}10$ configurations run at the training forcing on this skill--climate plane.

\textbf{H3 (an MLP with the same local input is not better than the polynomial): supported.} Offline $R^2$ is 0.850 against 0.847 (gap below the registered 0.01 threshold, although a paired $t$ test would detect it). At $c{=}10$ valid time is 1.421 against 1.398 ($p=$\rhval{cmp/main/mlp-32x2/f20_c10/valid_time/welch_p}) and spec\_err 0.137 against 0.144 ($p=$\rhval{cmp/main/mlp-32x2/f20_c10/spec_err/welch_p}); at $c{=}4$ the differences are even smaller (1.099 against 1.097, $p=$\rhval{cmp/main/mlp-32x2/f20_c4/valid_time/welch_p}). With 5 seeds this is absence of evidence of a difference rather than proof of equivalence; the paired test on valid time at $c{=}10$ is borderline ($p=$\rhval{cmp/main/mlp-32x2/f20_c10/valid_time/paired_p}). A one-input function of $X_k$ may simply be a smooth curve that a degree-4 polynomial already fits; we did not isolate this.

\textbf{H4 (stability, and MLP degradation under forcing shift): stability supported but uninformative; the shift part is refuted.} stable\_frac is 1.000 for every system, task and seed (Table~\ref{tab:main}) and in all ablation runs, so no closure diverged in 20 runs of 400 time units each; this does not exclude instabilities in other settings, in longer rollouts or with different inputs. Instabilities of offline-trained parametrisations have been reported for Earth system models~\cite{rasp2019coupled}; we saw none in this low-dimensional system, which says nothing about those models. Under a forcing shift (trained at $F{=}20$, deployed at $F{=}18$ and $22$; Table~\ref{tab:shift}) the MLP's w1\_pdf is lower than the polynomial's in both shifts (0.057 against 0.070 at 18, 0.034 against 0.039 at 22), the opposite of the hypothesis. The shift was mild: the deterministic closures' climates at $F{=}18$ and 22 are closer to truth than at the training value $F{=}20$ (e.g., polynomial w1\_pdf 0.070, 0.039 against \rhval{main/polynomial-deg-4/f20_c10/w1_pdf/mean} in the main group), so $F{=}20$ appears to be the hardest of the three regimes for local closures and this experiment does not test extrapolation properly.
